% section: 二項定理

  \section{二項定理}\label{節：二項定理}
  \noindent\textbf{【本編とのつながり】}\par
  和の累乗を展開する公式と組合せの意味を確認する.\sectionref*{節：極形式と複素平面}の\bookterm{euler}{オイラーの公式}の導出や,\sectionref*{節：発展難易度}の数え上げで使う.

    次に,和の累乗を展開するための公式を確認する.
    非負の整数$\,n$と実数$\,x,y$について,
    \[
      (x+y)^n=\sum_{k=0}^{n}\binom{n}{k}x^{n-k}y^{k}=\sum_{k=0}^{n}{}_n\mathrm{C}_kx^{n-k}y^{k}
    \]
    が成り立つ.これを\bookterm{binomial}{二項定理}という.\footnote{この多項式の式では,0乗は因子を含まない積を表すので,底が$0$でも$1$と読む.}
    ここで
\BookMathLayout{\[
      {}_n\mathrm{C}_k=\binom{n}{k}
      =\frac{n!}{k!(n-k)!}
      =\frac{n(n-1)\cdots(n-k+1)}{k!}
    \]}{
\[
\begin{aligned}
{}_n\mathrm C_k&=\binom nk=\frac{n!}{k!(n-k)!}\\
&=\frac{n(n-1)\cdots(n-k+1)}{k!}
\end{aligned}
\]
}
    が成り立つ.\,$\binom{n}{k}$および${}_n\mathrm{C}_k$は$\,n$個から$\,k$個を選ぶ組合せの数を表す.
    この階乗表示を使うのは,整数$k$が$0\le k\le n$を満たすときである.\,$0!=1$とする.
    $k=0$では,一つも選ばない選び方が1通りあるので$\binom n0=1$となる.
    また,本書では整数$k$が$k<0$または$k>n$を満たすとき,
    \[
      \binom nk=0
    \]
    と定める.その個数を選ぶ方法がないことに対応している.
    この場合に,階乗表示へ負の整数を代入してはいけない.
    例えば$\binom12=0$であり,
    \[
      \binom n2=\frac{n(n-1)}2\qquad(n\ge0)
    \]
    は$n=0,1$でも使える.
    \booknavref{節：総和型}{総和型}で現れた$n(n-1)/2$も,この書き方を使っている.
    \booknavref{節：二項係数と差分}{二項係数と差分}では,差を取ると添字が一つ下がることを確かめる.

    特に,\bookterm{binomial}{二項定理}ではこの組合せの数を二項係数と呼ぶ.
    その理由を,因子を選ぶ方法から考える.\,$(x+y)^n$は,\,$\,n$個の因子
    \[
      (x+y)(x+y)\cdots(x+y)
    \]
    の積である.展開するとき,各因子から$x$か$y$のどちらか一方を選ぶ.
    このとき,\,$\,x^{n-k}y^k$を作るには,\,$\,n$個の因子のうち$\,k$個から$y$を選び,
    残りの$\,n-k$個から$x$を選べばよい.
    $y$を選ぶ因子の位置の選び方は,\,$\,n$個から$\,k$個を選ぶ組合せの数
    $\binom{n}{k}$通りである.したがって,\,$\,x^{n-k}y^k$の係数は$\binom{n}{k}$となる.
    例えば$\,(x+y)^3$では,\,$\,x+y$という3個の因子のうち,
    どの2個から$x$を選ぶかによって$x^2y$が生じる.
    その選び方は$\binom{3}{2}=3$通りなので,\,$\,x^2y$の係数は$3$である.
    同様に,\,$\,xy^2$の係数も$\binom{3}{1}=3$となり,
    \[
      (x+y)^3=x^3+3x^2y+3xy^2+y^3
    \]
    を得る.
    また,これはパスカルの三角形の3段目の係数と一致する.
    複素数に対しても,同じ\bookterm{binomial}{二項定理}が成り立つ.
